\documentclass[../readings.tex]{subfiles}

\begin{document}

\subsection{Jacobi Eigenvalue Algorithm}
\label{sub:jacobi-theory}

Iterative algorithm for symmetric matrices using plane rotations to zero off-diagonal entries.

\subsubsection{Givens Rotations}

\addDef{Givens Rotation}{%
$G(p, q, \theta)$ is an identity matrix with a $2 \times 2$ rotation block at $(p, p), (p, q), (q, p), (q, q)$.
\begin{itemize}
    \item \textbf{Similarity Transform:}\enspace $\bA' = G^T \bA G$.
    \item \textbf{Effect:}\enspace Modifies only rows/columns $p$ and $q$. Preserves eigenvalues ($\bA'$ is orthogonally similar to $\bA$).
\end{itemize}
}
\label{def:givens-rotation}

\addThm{Optimal Rotation Angle}{%
To set $A'_{pq} = 0$, let $\tau = (A_{qq} - A_{pp}) / (2A_{pq})$. The optimal tangent $t = \tan\theta$ is:
\begin{align*}
    t = \frac{\text{sign}(\tau)}{|\tau| + \sqrt{1 + \tau^2}}, \quad c = \frac{1}{\sqrt{1+t^2}}, \quad s = ct.
\end{align*}
Choosing the smaller root $|t| \leq 1$ ensures numerical stability and minimal perturbation.
}
\label{thm:jacobi-angle}

\subsubsection{The Jacobi Algorithm}

\textBox{%
The Jacobi eigenvalue algorithm repeatedly applies orthogonal similarity transformations
\begin{align*}
    \bA^{(k+1)} = \bG_k^T \bA^{(k)} \bG_k
\end{align*}
to a symmetric matrix. Each rotation zeros one off-diagonal entry while preserving eigenvalues. As the off-diagonal entries converge to zero, $\bA^{(k)}$ converges to a diagonal matrix whose diagonal entries are the eigenvalues.
}

\addDef{Jacobi Eigenvalue Algorithm}{%
\textbf{Input:}\enspace Symmetric matrix $\bA\in\fR^{n\times n}$ and tolerance \texttt{tol}.\\
\textbf{Initialize:}\enspace $\bA^{(0)}=\bA$, $\bV^{(0)}=\bI$.\\
\textbf{Repeat:}
\begin{enumerate}
    \item Choose an off-diagonal pivot $(p,q)$ with $p<q$.
    \item Compute $c=\cos\theta$ and $s=\sin\theta$ using the stable tangent formula so that $(\bG^T\bA^{(k)}\bG)_{pq}=0$.
    \item Update the matrix by orthogonal similarity:
    \begin{align*}
        \bA^{(k+1)}=\bG^T\bA^{(k)}\bG.
    \end{align*}
    \item Accumulate eigenvectors:
    \begin{align*}
        \bV^{(k+1)}=\bV^{(k)}\bG.
    \end{align*}
\end{enumerate}
\textbf{Stop:}\enspace when $\sigma_{\text{off}}(\bA^{(k)}) \leq \texttt{tol}\|\bA^{(k)}\|_F$.\\
\textbf{Output:}\enspace Diagonal entries of $\bA^{(k)}$ approximate eigenvalues; columns of $\bV^{(k)}$ approximate eigenvectors.
}
\label{alg:jacobi-eigenvalue}

\addRemark{%
\textbf{(Pivot choices)}\enspace
\begin{itemize}
    \item \textbf{Classical Jacobi:}\enspace choose the largest off-diagonal entry $|A_{pq}|$.
    \item \textbf{Cyclic Jacobi:}\enspace sweep through all pairs $(p,q)$ in a fixed order.
\end{itemize}
The largest-entry rule gives the clearest convergence proof. Cyclic sweeps are simpler to implement and often used in practice.
}

\addRemark{%
\textbf{(Why eigenvectors accumulate)}\enspace After $k$ rotations,
\begin{align*}
    \bA^{(k)} = (\bV^{(k)})^T \bA \bV^{(k)}.
\end{align*}
If $\bA^{(k)}$ is nearly diagonal, then $\bA \bV^{(k)} \approx \bV^{(k)}\boldsymbol{\Lambda}$, so the columns of $\bV^{(k)}$ are approximate eigenvectors of the original matrix.
}

\subsubsection{Convergence and Cost}

\addThm{Monotone Convergence}{%
Let $\sigma^2_{\text{off}}(\bA) = \sum_{p < q} A_{pq}^2$ be the Frobenius norm of the off-diagonal entries. Each rotation $G(p, q, \theta)$ reduces this norm:
\begin{align*}
    \sigma^2_{\text{off}}(\bA') = \sigma^2_{\text{off}}(\bA) - A_{pq}^2.
\end{align*}
\textbf{Convergence:}\enspace Off-diagonals decrease monotonically. Convergence is ultimately \textbf{quadratic} (errors square each sweep).
}
\label{thm:jacobi-convergence}

\addPrf{%
The orthogonal transform $\bA' = G^T\bA G$ preserves the Frobenius norm: $\|\bA'\|_F^2 = \|\bA\|_F^2$. Let $d(\bA) = \sum A_{ii}^2$ be the sum of squares of diagonal entries. Then:
\begin{align*}
    \|\bA\|_F^2 = d(\bA) + 2\sigma_{\text{off}}^2(\bA).
\end{align*}
A Jacobi rotation $G(p,q,\theta)$ only affects rows and columns $p$ and $q$. For the $2 \times 2$ submatrix at $\{p,q\}$, the transformation is an eigendecomposition, so:
\begin{align*}
    (A'_{pp})^2 + (A'_{qq})^2 &= A_{pp}^2 + A_{qq}^2 + 2A_{pq}^2 \\
    d(\bA') &= d(\bA) + 2A_{pq}^2.
\end{align*}
Substituting this into the norm conservation $\|\bA'\|_F^2 = \|\bA\|_F^2$:
\begin{align*}
    d(\bA') + 2\sigma_{\text{off}}^2(\bA') &= d(\bA) + 2\sigma_{\text{off}}^2(\bA) \\
    (d(\bA) + 2A_{pq}^2) + 2\sigma_{\text{off}}^2(\bA') &= d(\bA) + 2\sigma_{\text{off}}^2(\bA) \\
    \sigma_{\text{off}}^2(\bA') &= \sigma_{\text{off}}^2(\bA) - A_{pq}^2.
\end{align*}
Thus each rotation reduces the off-diagonal mass by exactly $A_{pq}^2$.
}

\addRemark{%
\textbf{Performance vs.\ Accuracy:}\enspace
\begin{itemize}
    \item \textbf{Cost:}\enspace $O(n^3)$ per sweep. Typically 5--15 sweeps for convergence. More expensive than QR.
    \item \textbf{Accuracy:}\enspace Unlike QR, Jacobi achieves \textbf{high relative accuracy} on small eigenvalues.
\end{itemize}
}

\addExcr{%
\begin{enumerate}
    \item Apply one Jacobi rotation to zero $A_{12}$ in $\bA = \begin{pmatrix}2 & 1 \\ 1 & 3\end{pmatrix}$.
    \item Show that after zeroing $A_{pq}$, the new diagonal entries $A'_{pp}$ and $A'_{qq}$ are the eigenvalues of the $2 \times 2$ submatrix at $\{p, q\}$.
    \item Perform one full cyclic sweep on $\bA = \begin{pmatrix}1&1&0\\1&2&1\\0&1&3\end{pmatrix}$ by hand.
\end{enumerate}
}

\end{document}
